\documentclass[10pt,a4paper]{amsart}
\usepackage[margin=19mm]{geometry}
\usepackage{amsmath,amssymb,mathtools}
\usepackage[scr=rsfs,cal=euler]{mathalfa}
\usepackage{xfrac}
\usepackage[unicode,hidelinks,bookmarks=false]{hyperref}
\newcommand{\ie}{i.e.\ }
\newcommand{\cf}{cf.\ }
\newcommand{\resp}{respectively\ }
\newcommand{\Subsection}{\subsection}
\providecommand{\nobreakdash}{\nobreak\hbox{-}}
\setlength{\emergencystretch}{2em}
\multlinegap0pt
\usepackage{amssymb}
\usepackage{mathtools}
\usepackage[normalem]{ulem}
\usepackage{xcolor}
\usepackage{tikz}
\usepackage{cancel}
\usepackage[shortlabels]{enumitem}
\usepackage{float}
\newtheorem{lemma}{Lemma}
\newtheorem{corollary}{Corollary}
\newtheorem{theorem}{Theorem}
\newcommand{\calD}{\mathcal{D}}
\newcommand{\dist}{\text{dist}}
\newcommand{\hei}{\mbox{height}}
\newcommand{\hh}{\mbox{height}}
\newcommand{\pth}{\mbox{path}}
\newcommand\set[1]{\left\lbrace #1 \right\rbrace}
\title{Characterization of Well-Totally Dominated Trees}
\author{Jounglag Lim, James Gossell, Keri Ann Sather-Wagstaff, Devin Adams, Suzanna Castro-Tarabulsi, Aayahna Herbert, Vi Anh Nguyen, Yifan Qian, Matthew Schaller, Zoe Zhou, Yuyang Zhuo}
\date{Source: arXiv:2602.13956v1 (CC BY 4.0)}
\begin{document}
\maketitle

\noindent\emph{Source and licence.} Characterization of Well-Totally Dominated Trees. Version arXiv:2602.13956v1 (CC BY 4.0), distributed by arXiv under the Creative Commons Attribution 4.0 International licence, http://creativecommons.org/licenses/by/4.0/, as recorded in the arXiv licence metadata for this exact version and shown on the abstract page https://arxiv.org/abs/2602.13956v1. Full licence notice: \texttt{licenses/arxiv-2602.13956-CC-BY-4.0.txt}.\par\smallskip

\noindent\textit{Editorial scope.} Counted unit: the article's theorem WTD (source0.tex lines 898--902). The article's complete original proof of it is delivered with it (source0.tex lines 904--1036, one \texttt{proof} environment of 133 lines). No mathematics is written, repaired or re-derived: the statement and the proof are the article's own lines, in the article's own notation, and no retained line is substituted or re-flowed.  Retained uncounted as the source-local context the proof needs: lines 415--418; lines 596--598; lines 660--662; lines 746--749; lines 775--779; lines 880--884.  One declared adaptation: exactly 2 ancillary strings inside the retained spans are omitted (image-inclusion commands and, where the source repeats a label, the repeated label definitions) and no other line differs from the source; the figure environments, with their placement, captions and labels, are kept, and the package records the omitted strings in fidelity receipts bound to the affected bodies.  No retained line contains a citation, so the wrapper carries no bibliography and no bibliography entry.  The retained lines contain the article's own diagram code, which the wrapper typesets with the standard tikz packages; no image file is delivered and the package declares no asset dependency.  Editorial formatting only, in addition: an AMS \texttt{amsart} preamble that restates the article's own declarations and macros used by the retained lines, namely the environments \texttt{lemma}, \texttt{corollary}, \texttt{theorem} on separate counters, together with the standard packages those lines need.  The retained environments print in this document as Lemma 1, Corollary 1, Theorem 1 in order of appearance, whereas the article prints the same items with the numbering of its own full document.  The complete native source of the article as distributed in the arXiv e-print for this exact version is bundled unchanged as \texttt{sources/194672/source0.tex}.
\begin{lemma}\label{lem. D(v) not in intersection -> union is minimal}
    Let $D', D'' \subseteq V$ be disjoint minimal sets, and let $I = N(D') \cap N(D'')$.
    If there are domination selectors $\mathcal{D}':D' \to N(D)$ and $\mathcal{D}'': D'' \to N(D'')$ such that $\mathcal{D}'(D') \cap I = \emptyset$ and $\mathcal{D}''(D'') \cap I = \emptyset$, then $D = D' \cup D''$ is also minimal.
\end{lemma}

\begin{corollary}\label{cor. WTD iff RDBD WTD}
    A tree $T$ is WTD if and only if every minimal RDS has the same size and every minimal BDS has the same size.
\end{corollary}

\begin{theorem}\label{thm. P5 -> Mix}
Let $T$ be a balanced tree. If $T$ has two leaves $l_1$ and $l_2$ such that $\dist(l_1,l_2) = 4$, then $T$ is not WTD. 
\end{theorem}

\begin{lemma}\label{lem. no 4-path iff |N(v) cap V_1| = 1}
    Let $T$ be a balanced tree.
    Then for all leaves $\ell_1,\ell_2 \in V_0$, we have $\dist(\ell_1,\ell_2) \neq 4$ if and only if $|N(v) \cap V_1| = 1$ for all $v \in V_2$.
\end{lemma}

\begin{theorem}
\label{thm. height >= 4 is not WTD}
Let $T$ be a balanced tree with blue leaves and $\hh(T) \geq 4$. 
Then $T$ has two minimal RDSs of different sizes, and $T$ is not WTD.
\end{theorem}

\begin{theorem}
\label{thm. height(T)=2 -> not WTD}
Let $T$ be a balanced tree with $\hei(T) = 2$.
Then $T$ has two leaves of distance 4 apart, hence $T$ is not WTD.
\end{theorem}

\begin{theorem}
\label{thm. |N(s) cap V_2| > 1 -> not WTD}
Let $T$ be a balanced tree with blue leaves.
If there is a support vertex $s \in V_1$ such that $|N(s) \cap V_2| \geq 2$, then $T$ has two minimal RDSs of different sizes, hence $T$ is not WTD.
\end{theorem}

\begin{proof}
    If $\hh(T) \geq 4$, then $T$ has two minimal RDSs of different sizes by Theorem~\ref{thm. height >= 4 is not WTD}.
    So we assume that $\hh(T) \leq 3$.
    Since $V_2 \neq \emptyset$ by assumption, $\hh(T) \geq 2$.
    If $\hh(T) = 2$, then $T$ has two minimal RDSs of different sizes and is not WTD by Corollary~\ref{cor. WTD iff RDBD WTD} and Thoerem~\ref{thm. height(T)=2 -> not WTD}.
    Hence we assume that $\hh(T) = 3$.
    
    If there is a vertex $x \in V_2$ with $|N(x) \cap V_1| \geq 2$, then $T$ has two minimal RDSs of different sizes and is not WTD by Corollary~\ref{cor. WTD iff RDBD WTD}, Theorem~\ref{thm. P5 -> Mix}, and Lemma~\ref{lem. no 4-path iff |N(v) cap V_1| = 1}.
    Thus we also assume that for all $x \in V_2$, we have $|N(x) \cap V_1| = 1$.

    Let $l$ be a leaf adjacent to $s$.
    Since $|N(s) \cap V_2| \geq 2$, there are two distinct vertices $v_1,v_2$ in $N(s) \cap V_2$.
    Since $\hh(v_1) = \hh(v_2) = 2$, $v_1$ and $v_2$ are not leaves.
    Hence $\deg(v_1)$ and $\deg(v_2)$ are at least 2.
    By assumption, we have $|N(v_i) \cap V_1| = 1$ for $i = 1,2$.
    Thus $v_i$ must be adjacent to a height 3 vertex $u_i$ for $i= 1,2$.
    Similarly, since $u_1$ and $u_2$ are not leaves, there are distinct vertices $v_1',v_2' \in V_2$, distinct support vertices $s_1,s_2 \in V_1$, and distinct leaves $l_1,l_2 \in V_0$ such that $\pth(u_i,l_i) = (u_i,v_i',s_i,l_i)$ for $i = 1,2$.
    Let $T'$ be the subgraph of $T$ induced by the set $\set{l,l_1,l_2,s,s_1,s_2,v_1,v_1',v_2,v_2',u_1,u_2}$  (Figure~\ref{Figure3}).
    \begin{figure}[ht]
    \centering
    
    \caption{Induced subgraph $T'$ of $T$}
    \label{Figure3}
    \end{figure}
    
    Next, we construct a set dominating every red (odd height) vertex that is not in $T'$.
    Set
    \begin{align*}
        &X_3 := (N_T(V_2(T')) \cap V_3(T)) \setminus V_3(T') & & X_2 := (N_T(X_3) \cap V_2(T)) \setminus V_2(T')\\
        &X_1 := (N_T(X_2) \cap V_1(T)) \setminus V_1(T') & & X_0 := (N_T(X_1) \cap V_0(T))\setminus V_0(T')\\ 
        &Y_2 := (N_T(X_1) \cap V_2(T)) \setminus (X_2 \cup V_2(T')) &&
    \end{align*}
 see Figure~\ref{fig. support adj to V2 1}. 
 We claim that the set 
    $$
    U := (V_0(T) \cup V_2(T)) \setminus (X_0 \cup Y_2 \cup N_T(V_1(T') \cup V_3(T')))
    $$ 
    has $N_T(U) = (V_1(T) \cup V_3(T)) \setminus (V_1(T') \cup V_3(T'))$.
    \begin{figure}[ht]
    \centering
    
    \caption{Vertex sets $X_3$, $X_2$, $X_1$, $X_0$, and $Y_2$}
    \label{fig. support adj to V2 1}
    (Arrows indicate the adjacency relation between the sets; for instance, $Y_2 \subseteq N_T(X_1)$)
    \end{figure}
    
    The containment $N_T(U) \subseteq (V_1(T) \cup V_3(T)) \setminus (V_1(T') \cup V_3(T'))$ is straightforward since $U \subseteq V_0(T) \cup V_2(T)$ implies $N_T(U) \subseteq V_1(T) \cup V_3(T)$ because $T$ is a balanced tree, and $N_T(V_1(T') \cup V_3(T')) \cap U = \emptyset$ implies $N_T(U) \cap (V_1(T') \cup V_3(T')) = \emptyset$.

    Before we prove the reverse containment, we show that $X_2 \subseteq U$.
    Since $X_2 \subseteq V_2(T)$, we need to show that $X_2 \cap (X_0 \cup Y_2 \cup N_T(V_1(T') \cup V_3(T'))) = \emptyset$.
    By the constructions of $X_0$ and $Y_2$, we get $X_2 \cap (X_0 \cup Y_2) = \emptyset$.
    To show that $X_2 \cap N_T(V_1(T') \cup V_3(T')) = \emptyset$, assume by way of contradiction that there is a vertex $c \in X_2 \cap N_T(V_1(T') \cup V_3(T'))$.
    We construct a cycle containing the vertex $c$, contradicting that $T$ is a tree.
    First, suppose that $c \in N_T(V_1(T'))$.
    Then there is a vertex $s' \in V_1(T')$ that is adjacent to $c$.
    By construction of $X_2$ and $X_3$, there are vertices $x_3 \in X_3$ and $v'$ in $V_2(T')$ such that $cx_3 \in E(T)$ and $v'x_3 \in E(T)$.
    Since $T'$ is a tree and $s'$ and $v'$ are vertices in $T'$, there is a path $P$ in $T'$ from $v'$ to $s'$.
    Thus the paths $P$ and $(s',c,x_3,v')$ form a cycle in $T$.
    Similarly, we can construct a cycle when $c \in N_T(V_3(T'))$.

    Now we show that $N_T(U) \supseteq (V_1(T) \cup V_3(T)) \setminus (V_1(T') \cup V_3(T'))$.
    Let $x \in (V_1(T) \cup V_3(T))$ such that $x \not\in (V_1(T') \cup V_3(T'))$.
    We show that there is a vertex in $U$ that is adjacent to $x$.
    We have three cases to consider.

    \emph{Case 1.} $x \in X_3 \cup X_1$.
    If $x \in X_1$, then there is a vertex $x_2 \in X_2 \subseteq U$ that is adjacent to $x$.
    If $x \in X_3$, then $\deg(x) \geq 2$ as $x$ is not a leaf.
    Hence $x$ is adjacent to at least 2 vertices in $V_2(T)$.
    By construction of $X_3$, $x$ is adjacent to a vertex in $V_2(T')$.
    Notice that $x$ cannot be adjacent to 2 vertices in $V_2(T')$ since this will form a cycle in $T$.
    Thus $x$ must be adjacent to a vertex $y \in V_2(T)$ such that $y \not\in V_2(T')$.
    By construction of $X_2$, $y \in X_2 \subseteq U$ as desired.

    \emph{Case 2.} $x \in V_1(T)\setminus X_1$. Then there is a leaf $l_x \in V_1(T)$ that is adjacent to $x$.
    Since $x \not\in X_1$, $l_x \not\in X_0$ as $x$ is the only vertex $l_x$ is adjacent to.
    We also have $l_x \not\in Y_2 \cup N_T(V_3(T'))$ since vertices in $Y_2 \cup N_T(V_3(T'))$ have height 2.
    Finally, $l_x \not\in N_T(V_1(T'))$ since we assumed that $x \not\in V_1(T')$.
    Thus we have $l_x \in U$ as desired.
    
    \emph{Case 3.} $x \in V_3(T)\setminus X_3$.
    Then $x$ is not a leaf, hence $\deg_T(x) \geq 2$. 
    Also, as $\hh(T) = 3=\hh(x) = 3$ and $T$ is balanced, $x$ is adjacent to at least two height 2 vertices $t_1,t_2 \in V_2(T)$.
    To show that at least one of $t_1$ or $t_2$ is in $U$, assume by way of contradiction that $t_1,t_2 \not\in U$.
    Then $t_1,t_2 \in X_0 \cup Y_2 \cup N_T(V_1(T') \cup V_3(T'))$ by definition of $U$.
    Since height of $t_1$ and $t_2$ is 2, we have $t_1,t_2 \in Y_2 \cup N_T(V_1(T') \cup V_3(T'))$.
    As in the proof of $X_2 \subseteq U$ there is a path from $t_1$ to $t_2$ only using the vertices in $V(T') \cup \left(\bigcup_{i = 1}^3X_i\right) \cup Y_2 \cup N_T(V_1(T') \cup V_3(T'))$. 
    But $(t_1, x, t_2)$ is also a path from $t_1$ to $t_2$, with $x$ not belonging to $V(T') \cup \left(\bigcup_{i = 1}^3X_i\right) \cup Y_2 \cup N_T(V_1(T') \cup V_3(T'))$.
    Thus we have a cycle in $T$, a contradiction.

    Therefore, we have $N_T(U) = (V_1(T) \cup V_3(T)) \setminus (V_1(T') \cup V_3(T'))$ as desired.

    
    Now let $D \subseteq U$ be a minimal set of $U$.
    Let $D_1 := \set{l_1,l_2,v_1,v_2}$ and $D_2 := \set{l,v_1',v_2'}$.
    We claim that the sets $D \cup D_1$ and $D \cup D_2$ are minimal RDSs.
    For $i = 1,2$, since $T$ is balanced and $N_T(D_i) \supseteq \set{u_1,u_2,s,s_1,s_2} = V_1(T') \cup V_3(T')$, we get $N_T(D \cup D_i) = V_1(T) \cup V_3(T)$.

    To show the minimality of $D \cup D_1$ and $D \cup D_2$, define domination selectors for $D_1$ and $D_2$
    $$
    \calD_1(v) := \begin{cases}
    s_1 & v = l_1\\
    s_2 & v = l_2\\
    u_1 & v = v_1\\
    u_2 & v = v_2
    \end{cases}
    \qquad
    \calD_2(v) := \begin{cases}
        s_1 & v = v_1'\\
        s_2 & v = v_2'\\
        s & v = l\ \ \ .
    \end{cases}
    $$

    For $D$, we claim that there is a domination selector $\calD: D \to N_T(D)$ such that for all $x \in D \cap X_2$, we have $\calD(x) \in X_1$.
    To this end, we need to show that for all $x \in D \cap X_2$, there is a vertex $y \in X_1$ such that $N_T(y) \cap D = \set{x}$.
    So, assume by way of contradiction that there is a vertex $x \in D \cap X_2$ such that for all $y\in X_1$, we have $N_T(y) \cap D \neq \set{x}$.
    Since $\hh(x) = 2$, there is a support vertex $s' \in V_1(T)$ that is adjacent to $x$.
    By construction of $X_2$, either $s' \in X_1$ or $s' \in V_1(T')$.
    To show that $s' \in X_1$, assume by way of contradiction that $s' \in V_1(T)'$.
    Then one can show that this assumption creates a cycle containing $s'$ and $x$ using the vertices in $V(T') \cup X_2 \cup X_3$, a contradiction.
    Thus $s' \in X_1$.
    Since $N_T(s') \cap D \neq \set{x}$ but $x \in N_T(s') \cap D$, there is a vertex $x' \in D \cap X_2$ such that $N_T(s') \cap D \supseteq \set{x,x'}$.
    By the construction of $X_2$ and $X_3$, there is a path $P$ from $x$ to $x'$ only using the vertices in $V(T') \cup X_3$.
    On the other hand, $(x,s',x')$ is a path from $x$ to $x'$ which is not $P$ as $s' \not \in V(T') \cup X_3$.
    Hence we have a cycle in $T$, a contradiction, and the claim is established.

    Furthermore, since the only vertices in $U$ that are adjacent to $X_3$ are in $X_2$, we get $\calD(D) \cap X_3 = \emptyset$ as $\calD(x) \in X_1$ for all $x \in D \cap X_2$.

    Now consider $N_T(D) \cap D_i$ for $i = 1,2$.
    By the above computation of $N_T(U) = N_T(D)$, we have $N_T(D) \cap D_i =\emptyset$ for $i = 1,2$.
    Since $\calD_i(D_i) \cap X_3 = \emptyset$ for $i = 1,2$, and $\calD(D) \cap X_3 = \emptyset$, the sets $D \cup D_1$ and $D \cup D_2$ are minimal by Lemma~\ref{lem. D(v) not in intersection -> union is minimal}. 
\end{proof}  

\end{document}
